\section{Theoretical Framework}
\label{sec:theory}

\subsection{The Paired Efficiency Factor}

\subsubsection{Derivation}

For random variables $X_{\mathrm{A}}$ and $X_{\mathrm{B}}$ with variances $\sigma^{2}_{\mathrm{A}}$, $\sigma^{2}_{\mathrm{B}}$ and correlation $\rho$,
\begin{equation}
  \Var(X_{\mathrm{A}} - X_{\mathrm{B}}) = \sigma^{2}_{\mathrm{A}} + \sigma^{2}_{\mathrm{B}} - 2\rho\sigma_{\mathrm{A}}\sigma_{\mathrm{B}}.
  \label{eq:var_sum}
\end{equation}
Introducing $\kappa=\sigma^{2}_{\mathrm{B}}/\sigma^{2}_{\mathrm{A}}$ and factoring,
\begin{equation}
  \Var(X_{\mathrm{A}} - X_{\mathrm{B}}) = \sigma^{2}_{\mathrm{A}}\bigl(1 + \kappa - 2\sqrt{\kappa}\,\rho\bigr).
  \label{eq:var_factored}
\end{equation}

\noindent Here $\sigma^{2}_{\mathrm{A}}$ is the variance of entity~A's measurements, $\kappa=\sigma^{2}_{\mathrm{B}}/\sigma^{2}_{\mathrm{A}}$ is the variance ratio, and $\rho$ is the Pearson correlation between $X_{\mathrm{A}}$ and $X_{\mathrm{B}}$. For unpaired measurements ($\rho=0$), the variance is $\sigma^{2}_{\mathrm{A}}(1+\kappa)$. The Paired Efficiency Factor is the ratio of unpaired to paired variance, i.e.\ the expression in~\cref{eq:pef}:
\begin{equation*}
  \eta = \frac{1+\kappa}{1+\kappa-2\sqrt{\kappa}\,\rho}.
\end{equation*}
\noindent with $\eta$, $\kappa$, and $\rho$ as defined for \cref{eq:pef} in \cref{sec:pef_def}.

\medskip
\noindent\textbf{Remark.} No distributional assumptions are required: \cref{eq:pef,eq:var_factored} hold for any random variables with finite second moments.

\subsubsection{Properties}

PEF is strictly increasing in $\rho$ for fixed $\kappa>0$,
\begin{equation}
  \frac{\partial \eta}{\partial \rho} = \frac{2\sqrt{\kappa}(1+\kappa)}{\bigl(1+\kappa-2\sqrt{\kappa}\,\rho\bigr)^{2}} > 0,
  \label{eq:detadrho}
\end{equation}
and convex in $\rho$,
\begin{equation}
  \frac{\partial^{2} \eta}{\partial \rho^{2}} = \frac{8\kappa(1+\kappa)}{\bigl(1+\kappa-2\sqrt{\kappa}\,\rho\bigr)^{3}} > 0.
  \label{eq:d2etadrho2}
\end{equation}
Three regimes arise:
\begin{itemize}
  \item $\eta>1$: pairing reduces variance (positive correlation dominates);
  \item $\eta=1$: pairing provides no efficiency change (Fisher baseline);
  \item $\eta<1$: pairing increases variance (negative correlation dominates).
\end{itemize}
\Cref{fig:pef_landscape} visualises the PEF landscape across the $(\rho,\kappa)$ parameter space.

\subsubsection{Reduction to Classical Case}

When $\kappa=1$,
\begin{equation*}
  \eta = \frac{2}{2-2\rho} = \frac{1}{1-\rho},
\end{equation*}
recovering Fisher's classical paired efficiency formula \parencite{fisher1935}.

\subsubsection{Provenance and Scope of Contribution}

The variance-of-difference identity for unequal variances (\cref{eq:var_sum,eq:var_factored}) is a standard second-moment result, and the efficiency-of-pairing argument under equal variances is due to Fisher \parencite{fisher1935}; the same paired-margin machinery underlies the Pitman--Morgan test for equality of correlated variances \parencite{pitman1939,morgan1939}. Our contribution is not this algebra but its reframing: we collect it into a single closed-form efficiency ratio (\cref{eq:pef}), parameterised by $(\kappa,\rho)$, and repurpose it from an inferential setting toward feature construction and predictive information in competitive measurement. Conceptualising PEF as a formalisation built on these seminal results, rather than as a new identity, is a deliberate strength of the framework. We do not claim priority for the underlying algebra.

Readers familiar with design efficiency will recognise \cref{eq:pef} as a variance ratio; the framework built around it---the $(\kappa,\rho)$ taxonomy, its link to predictive information, and the resulting feature-engineering guidance (enumerated in \cref{sec:contributions})---is the primary contribution of this paper. One further anchor specific to the variance-ratio formula is established here: under normality, $\eta$ coincides with the Pitman asymptotic relative efficiency of the paired versus independent $t$-test \parencite{lehmann1999} (\cref{sec:pitman}), a testing-theory characterisation absent from prior treatments.

\subsection{The PEF--Information Content Relationship}

\subsubsection{Setup and Assumptions}
\label{sec:mi_setup}

We now impose distributional assumptions to connect PEF with predictive power. Let $Y\in\{0,1\}$ denote the \textbf{binary match outcome} (whether entity~A outperforms entity~B on the measure of interest), and let $X=X_{\mathrm{A}}-X_{\mathrm{B}}$ be the corresponding relative feature. The predictive setting introduces a quantity beyond the second-moment pair $(\kappa,\rho)$ that governs PEF: the \textbf{signal}, or mean separation $\delta=\mu_{\mathrm{A}}-\mu_{\mathrm{B}}$ between the two entities. Whereas $(\kappa,\rho)$ fix the variance geometry of the relative feature, it is the magnitude of $\delta$ relative to $\sigma_{\mathrm{A}}$ that determines how much of that geometry translates into predictive power; this distinction is the crux of the efficiency--power tension and its consequences are developed in \cref{sec:signal_strength}. We require two assumptions:

\begin{description}
  \item[(A1)] $(X_{\mathrm{A}},X_{\mathrm{B}})$ follows a bivariate normal distribution with means $\mu_{\mathrm{A}}$, $\mu_{\mathrm{B}}$, variances $\sigma^{2}_{\mathrm{A}}$, $\sigma^{2}_{\mathrm{B}}$, and correlation $\rho$. Under (A1), $X=X_{\mathrm{A}}-X_{\mathrm{B}}$ is normally distributed with mean $\delta=\mu_{\mathrm{A}}-\mu_{\mathrm{B}}$ and variance given by \cref{eq:var_factored}.
  \item[(A2)] The class-conditional distributions of $X$ given $Y$ follow a symmetric Gaussian discriminant model:
  \[
    X \mid Y=1 \;\sim\; \mathcal{N}\!\left(+\tfrac{\delta}{2},\,\Var(X)\right),
    \quad
    X \mid Y=0 \;\sim\; \mathcal{N}\!\left(-\tfrac{\delta}{2},\,\Var(X)\right),
  \]
  with equal class priors $\mathbb{P}(Y=1)=\mathbb{P}(Y=0)=\tfrac{1}{2}$.
\end{description}

Under (A2), $Y$ is \emph{not} a deterministic function of $X$: the outcome depends on many factors beyond the single observed KPI difference, so $H(Y\mid X)>0$. The Gaussian discriminant model is the natural linear classifier for normally distributed features and is consistent with (A1) when the within-class variance equals the marginal $\Var(X)$.

\medskip
\noindent\textbf{Remark.} Assumption (A1) is empirically reasonable for the present data: the within-match KPI difference series $X=X_{\mathrm{A}}-X_{\mathrm{B}}$ in professional rugby and football are approximately normal (mean Shapiro--Wilk $W$ of $0.94$--$0.95$, near-zero skewness), consistent with the central-limit behaviour of paired differences even where the individual team series are mildly skewed (Supplementary Information). Note that (A1) is required only for the information content derivation; the PEF formula itself (\cref{eq:pef}) is distribution-free. Assumption (A2) is an additional approximation for the predictive setting and is likewise not required for the PEF formula.

\subsubsection{Derivation}

The mutual information \parencite{shannon1948,cover2006} between $X$ and $Y$ is
\begin{equation}
  I(X;Y) = H(Y) - H(Y\mid X).
  \label{eq:mi_def}
\end{equation}

\noindent Here $I(X;Y)$ is the mutual information (in bits) between the relative feature $X$ and the binary outcome $Y$; $H(Y)$ is the Shannon entropy of $Y$; and $H(Y\mid X)=\mathbb{E}_{X}[H(Y\mid X=x)]$ is the conditional entropy. For equiprobable outcomes, $H(Y)=1$ bit. Under (A2), the posterior probability of winning given the observed KPI difference is
\[
  \mathbb{P}(Y=1\mid X=x) = \Phi\!\left(\frac{x}{\sqrt{\Var(X)}}\right),
\]
where $\Phi$ is the standard normal CDF. The Bayes error rate is accordingly $\Phi\!\bigl(-\delta/(2\sqrt{\Var(X)})\bigr)$, and the expected conditional entropy evaluates to
\begin{equation}
  H(Y\mid X) = H\!\left(\Phi\!\left(\frac{\delta}{2\sqrt{\Var(X)}}\right)\right),
  \label{eq:hygx}
\end{equation}
where $H(p)=-p\log_{2}p-(1-p)\log_{2}(1-p)$ is the binary entropy function. (Full derivation in \cref{sec:appendix}.)

Substituting $\Var(X)=\sigma^{2}_{\mathrm{A}}(1+\kappa)/\eta$ from the rearranged PEF formula,
\begin{equation}
  I(X;Y) = 1 - H\!\left(\Phi\!\left(\frac{\delta}{2\sigma_{\mathrm{A}}\sqrt{(1+\kappa)/\eta}}\right)\right).
  \label{eq:mi_closed}
\end{equation}

\noindent Here $\delta$, $\sigma_{\mathrm{A}}$, $\Phi$, $H(p)$, $\kappa$, and $\eta$ retain their earlier definitions (\cref{sec:mi_setup}; \cref{eq:pef,eq:hygx}). This is the \textbf{PEF--information content relationship}. $\eta$ and $I(X;Y)$ are functionally linked through the signal-to-noise ratio $\delta/\bigl(2\sigma_{\mathrm{A}}\sqrt{(1+\kappa)/\eta}\bigr)$. \Cref{fig:info_surface} visualises the information content surface $I(X;Y)$ across the $(\kappa,\rho)$ parameter space for the nominal case $\delta/\sigma_{\mathrm{A}}=1$; sensitivity of the surface to this parameter is examined in the Supplementary Information (Figure~S1).

\subsubsection{The Efficiency--Power Tension}
\label{sec:eff_power_theory}

From \cref{eq:mi_closed}, $\eta$ enters the information content formula through the denominator $\sqrt{(1+\kappa)/\eta}$. Consider two effects of negative correlation ($\rho<0$):
\begin{enumerate}
  \item \textbf{Variance increases:} $\eta<1$, meaning $\Var(X)>\Var_{\text{unpaired}}(X)$. This harms statistical efficiency.
  \item \textbf{Signal may increase relative to noise:} if the mean difference $\delta$ is large relative to the amplified standard deviation, $I(X;Y)$ can still be substantial.
\end{enumerate}
Predictive power depends on $I(X;Y)$, not $\eta$ alone. When $\delta$ is sufficiently large relative to the amplified standard deviation, the information content remains high despite increased variance. This accounts for cases where machine learning models benefit from relative features even when $\eta<1$.

\subsection{Sources of Correlation in Paired Comparisons}

Paired measurements acquire correlation through three generic mechanisms, each of which appears across competitive and non-competitive domains:
\begin{itemize}
  \item \textbf{Shared exogenous factors.} Conditions common to both entities move their measurements together, inducing positive correlation. Examples include market-wide movements in finance, batch, instrument, or site effects in manufacturing and clinical measurement, and weather, venue, or officiating in sport.
  \item \textbf{Zero-sum interaction.} When one entity's gain is the other's loss, their measurements move oppositely, inducing negative correlation. Examples include rivals competing for a fixed market share, and head-to-head contests in which possession or territory ceded by one competitor accrues directly to the other.
  \item \textbf{Non-zero-sum (common-mode) constraints.} When both entities can succeed or fail together, positive correlation can arise even under competition. Examples include firms lifted or depressed by a common sector trend, and matches in which both teams perform well in open, high-tempo play.
\end{itemize}
The sign and magnitude of $\rho$ therefore reflect which mechanism dominates in a given setting. The PEF formula accommodates all correlation signs within a single expression, covering shared-factor regimes (typically $\rho>0$) and zero-sum competitive regimes (potentially $\rho<0$). The competitive sports KPIs analysed in \cref{sec:results} provide concrete instances spanning this range, including the negatively correlated Quadrant~4 cases that are the focus of this paper. In the primary sports application, each pair comprises the two teams' KPI values from the same fixture, so the estimated $\rho$ reflects within-match co-movement between competitors rather than serial dependence of a single team across games; estimation scope and the associated independence considerations are detailed in \cref{sec:methods}.

\subsection{Quadrant Taxonomy}

The $(\kappa,\rho)$ parameter space divides naturally into four quadrants with distinct statistical and predictive properties (\cref{tab:quadrants}):

\textbf{Quadrant 1 ($\kappa>1$, $\rho>0$): High efficiency, high information.} Positive correlation combined with variance asymmetry produces $\eta>1$ with high information content. Relative features outperform absolute features in this regime. Exemplified by market-adjusted returns and paired clinical trials.

\textbf{Quadrant 2 ($\kappa<1$, $\rho>0$): High efficiency, moderate information.} Positive correlation provides variance reduction even with low variance ratio, though information content is lower than Quadrant~1. Exemplified by manufacturing control charts.

\textbf{Quadrant 3 ($\kappa<1$, $\rho<0$): Low efficiency, low information.} Negative correlation with low variance asymmetry amplifies variance without commensurate information gain. Absolute features are preferred. Exemplified by weakly competitive or uncorrelated independent measurements.

\textbf{Quadrant 4 ($\kappa>1$, $\rho<0$): Low efficiency, variable information.} This is where the efficiency--power tension is most pronounced. $\eta<1$ indicates statistical harm, yet information content may remain substantial when $\delta$ is large. The decision to relativise requires explicit analysis of both $\eta$ and $I(X;Y)$. Exemplified by competitive sports and head-to-head business comparisons.

\subsection{Signal Strength and the Limits of \texorpdfstring{$\eta$}{eta} as a Predictor}
\label{sec:signal_strength}

The PEF--information content relationship (\cref{eq:mi_closed}) shows that $\eta$ is not the sole determinant of predictive gain. The mutual information $I(X;Y)$ depends on $\eta$ through the denominator $\sqrt{(1+\kappa)/\eta}$, but also on the signal-to-noise ratio
\begin{equation}
  \frac{\delta}{\sigma_{\mathrm{A}}} = \frac{\mu_{\mathrm{A}} - \mu_{\mathrm{B}}}{\sigma_{\mathrm{A}}},
  \label{eq:delta_ratio}
\end{equation}
the standardised mean difference between the two entities. Two KPIs with identical $(\kappa,\rho)$ positions---and therefore identical $\eta$---will yield different information content if their signal strengths differ.

This has a practical consequence: $\eta$ alone does not determine ML gain. $\eta$ encodes the $(\kappa,\rho)$ position and identifies the quadrant regime, but two KPIs in the same quadrant will yield different ML outcomes if their signal strengths differ. The appropriate diagnostic is therefore the $(\kappa,\rho,\delta/\sigma_{\mathrm{A}})$ triple: quadrant for the directional prediction, $\delta/\sigma_{\mathrm{A}}$ estimated from data for the magnitude. Supplementary~\cref{fig:si_iso_eta_I} shows how iso-$\eta$ and iso-$I(X;Y)$ contours diverge once $\delta/\sigma_{\mathrm{A}}$ varies.

The practical implication is that the four-quadrant taxonomy and the $I(X;Y)$ formula together constitute the appropriate predictive tool, with $\delta/\sigma_{\mathrm{A}}$ estimated from data. \Cref{sec:exemplars} illustrates this with four quadrant exemplars spanning the empirical range of $(\kappa,\rho,\delta/\sigma_{\mathrm{A}})$ values.

\subsection{Cross-Domain Connections}

Under the PEF framework, established practices across disciplines can be viewed as instances of a common principle:

\begin{center}
\begin{tabular}{@{}llll@{}}
\toprule
\textbf{Field} & \textbf{Method} & \textbf{Correlation source} & \textbf{Typical PEF regime} \\
\midrule
Statistics & Paired $t$-test & Repeated measures & $\eta=1/(1-\rho)$ when $\kappa=1$ \\
Finance & Market-adjusted returns & Market factors & Q1: $\eta>1$, high $\rho$ \\
Healthcare & Paired clinical trials & Temporal/biological & Q1--Q2: $\eta\gg 1$ \\
Manufacturing & Control charts & Process dynamics & Q2: $\eta>1$, moderate $\rho$ \\
Sports & Relative metrics & Environmental/competitive & Q4: $\eta$ variable \\
\bottomrule
\end{tabular}
\end{center}

When two domains share similar $(\kappa,\rho)$ values, techniques from one may transfer to the other, for example, methods from healthcare (high-$\rho$ settings) to manufacturing contexts with similar correlation structures.
